\documentclass{article}
\usepackage[utf8]{inputenc}
\usepackage{a4wide}

\begin{document}

\section*{เมล์ถึงเพื่อนรัก}

ลภิณรัตน์เขียน email ถึงเพื่อนชาวต่างชาติ แต่ด้วยความขี้เล่นและชอบแกล้งเพื่อนของเธอ

เธอจึงพยายามปรับแต่งข้อความใน email ให้ดูแปลกประหลาดมากขึ้น

โดยการจัดเรียงตัวอักษรในข้อความให้เป็นไปตามกฎที่กำหนดว่า

"ตัวอักษรพิมพ์ใหญ่ต้องอยู่ก่อนตัวอักษรพิมพ์เล็ก"

ซึ่งหมายความว่าในข้อความของเธอทั้งหมดต้องจัดเรียงอย่างถูกต้องตามกฎตัวอักษรพิมพ์ใหญ่จะต้องอยู่ในตำแหน่งก่อนตัวอักษรพิมพ์เล็กเสมอ



เพื่อให้ได้ข้อความตามกฎที่เธอตั้งไว้

ลภิณรัตน์สามารถทำการเปลี่ยนแปลงตัวอักษรจากพิมพ์ใหญ่เป็นพิมพ์เล็ก

หรือจากพิมพ์เล็กเป็นพิมพ์ใหญ่ในตำแหน่งใดก็ได้ภายในข้อความ

โดยการเปลี่ยนแปลงตัวอักษรต้องเป็นตัวอักษรเดิมเท่านั้น

นี่คือสิ่งเดียวที่เธอทำได้ในการแปลงข้อความให้เป็นไปตามกฎที่ตั้งไว้



\ul{งานของคุณ}:

ช่วยลภิณรัตน์หาจำนวนครั้งที่น้อยที่สุดที่เธอต้องเปลี่ยนแปลงตัวอักษรในข้อความเพื่อให้ได้ข้อความที่เป็นไปตามกฎที่กำหนด



\subsection*{Input}
บรรทัดเดียว - ข้อความความยาวไม่เกิน 10\^{}5 ตัวอักษรที่ไม่มีช่องว่าง

และประกอบด้วยตัวอักษรภาษาอังกฤษพิมพ์เล็กหรือพิมพ์ใหญ่เท่านั้น\\



\subsection*{Output}
บรรทัดเดียว -

จำนวนครั้งที่น้อยที่สุดที่เธอต้องเปลี่ยนแปลงตัวอักษรในข้อความเพื่อให้ได้ข้อความที่เป็นไปตามกฎที่กำหนด\\



\end{document}
